\documentclass[10pt,a4paper]{amsart}
\usepackage[margin=19mm]{geometry}
\usepackage{amsmath,amssymb,mathtools}
\usepackage[scr=rsfs,cal=euler]{mathalfa}
\usepackage{xfrac}
\usepackage[unicode,hidelinks,bookmarks=false]{hyperref}
\newcommand{\ie}{i.e.\ }
\newcommand{\cf}{cf.\ }
\newcommand{\resp}{respectively\ }
\newcommand{\Subsection}{\subsection}
\providecommand{\nobreakdash}{\nobreak\hbox{-}}
\setlength{\emergencystretch}{2em}
\multlinegap0pt
\usepackage{bm}
\usepackage{bbm}
\usepackage{dsfont}
\usepackage{xspace}
\usepackage{cancel}
\usepackage{mathtools}
\usepackage{amsthm}
\makeatletter
\@ifundefined{thm}{\newtheorem{thm}{Theorem}}{}
\@ifundefined{defi}{\newtheorem{defi}[thm]{Definition}}{}
\@ifundefined{remark}{\newtheorem{remark}[thm]{Remark}}{}
\makeatother
\def\xsum{\mathop{\sum\nolimits'}}
\newcommand{\N}{\mathbb N}
\newcommand{\R}{\mathbb R}
\title[Structural transitions in interacting lattice systems]{Structural transitions in interacting lattice systems}
\author{Laurent B\'etermin, Ladislav \v{S}amaj, Igor Trav\v{e}nec}
\date{Source: arXiv:2312.01395v2 (CC BY 4.0)}
\begin{document}
\maketitle

\noindent\emph{Source and licence.} Structural transitions in interacting lattice systems. Version arXiv:2312.01395v2 (CC BY 4.0), distributed under the Creative Commons Attribution 4.0 International licence, as recorded in the arXiv licence metadata for this exact version and shown on the abstract page https://arxiv.org/abs/2312.01395v2. Full rights notice: licenses/arxiv-2312.01395-CC-BY-4.0.txt.\par\smallskip

\noindent\textit{Editorial scope.} Counted unit: the Remark whose source label is \texttt{None} in the article (source0.tex lines 497--499, source label \texttt{None}), delivered together with the article's own complete original proof of it (source0.tex lines 500--530, one \texttt{proof} environment of 31 lines). No mathematics is written, repaired or re-derived: statement and proof are the article's own text line for line. Required context retained uncounted: ereca (lines 247--254); regularity (lines 260--262); ecreps (lines 350--385). Every definition, display and label of those lines is retained unchanged. Omitted from this package and not redistributed: 1--246, 255--259, 263--349, 386--496, 531--1164. Nothing inside a retained excerpt is omitted or altered: every retained body is delivered exactly as the article's lines are written, so no omission receipt of any kind accompanies this package. Citations: the retained excerpts carry no citation command at all, so no bibliography entry of the article is transcribed and no reference is redistributed. Reference adaptation: none was needed and none was made; every reference inside the delivered unit resolves inside it, so no unresolved reference or citation is produced. Printed slips kept literal: no printed word, symbol, hypothesis or number of the counted statement or of its proof was altered. Layout and class: the article's own class is \texttt{interact} with packages including amsmath, latexsym, amsthm, graphics, graphicx, subfigure, amssymb, adjustbox, a4wide, color; the wrapper is the lane's pinned \texttt{amsart} editorial wrapper at 10pt on A4, which restates the theorem-like environments the retained text needs and adds the article's own macro definitions that the retained text needs (\texttt{\textbackslash N}, \texttt{\textbackslash R}, \texttt{\textbackslash xsum}), with extra wrapper packages (bm, bbm, dsfont, xspace, cancel, mathtools). The delivered numbering is the unit's own. Assets: no retained span contains a figure environment, a graphics-inclusion command, a PDF-page inclusion, a table or a TikZ picture; no image file is copied into this package, the package declares no asset dependency and no image file is redistributed with it. Licence: version arXiv:2312.01395v2 of this article is distributed under the Creative Commons Attribution 4.0 International licence, as recorded in the arXiv licence metadata for this exact version and shown on the abstract page https://arxiv.org/abs/2312.01395v2. A full rights notice is delivered at licenses/arxiv-2312.01395-CC-BY-4.0.txt. Characters: every retained excerpt is pure ASCII, so the delivered engine, XeLaTeX with its default Latin Modern text font, reports no missing character.
\begin{defi}[\textbf{Interaction energy}]
Let $f\in \mathcal{F}$, $A>0$, $\Delta\in (0,1]$, then the $f$-energy of $L_{A,\Delta}$ is defined by
\begin{equation} \label{ereca}
E(A,\Delta):=\frac{1}{2} \xsum_{j,k=-\infty}^{\infty}
f\left(\sqrt{Q_{A,\Delta}(j,k)} \right)=\frac{1}{2} \xsum_{j,k=-\infty}^{\infty}
f \left( \sqrt{A \left( \frac{j^2}{\Delta} + k^2\Delta \right)} \right).
\end{equation}
\end{defi}

\begin{remark}\label{regularity}
Since $f$ is admissible, it follows immediatly that, by composition and absolute summability of the energy \eqref{ereca}, $E\in C^{\infty}(\R_+^*\times \R_+^*)$.
\end{remark}

\begin{thm}[\textbf{Taylor expansion of the energy}]
Let $f\in \mathcal{F}$ and $A>0$, then, as $\varepsilon\to 0$,
\begin{equation} \label{ecreps}
E(A,e^\varepsilon) = E_0(A) + E_2(A) \varepsilon^2 + E_4(A) \varepsilon^4
+ O(\varepsilon^6) ,
\end{equation}
where 
\begin{equation} 
E_0(A) = E(A,1)= \frac{1}{2} \int_0^\infty 
\left[ \left( \theta_3 \right)^2 - 1 \right]{\rm d}\mu_f\left(\frac{t}{A} \right)
\end{equation}
is the energy of the square lattice,
\begin{equation} \label{E2A} 
E_2(A) =\frac{d^2}{d\varepsilon^2}\left[ E(A,e^\varepsilon)\right]_{|\varepsilon=0}
= \frac{1}{2} \int_0^\infty \left[ t \theta_3 \theta_3^{(1)}
- t^2 \left( \theta_3^{(1)} \right)^2 + t^2 \theta_3 \theta_3^{(2)} \right]
{\rm d}\mu_f\left(\frac{t}{A} \right) ,
\end{equation}
and 
\begin{eqnarray} 
E_4(A)  = \frac{d^4}{d\varepsilon^4}\left[ E(A,e^\varepsilon)\right]_{|\varepsilon=0}
& = & \frac{1}{24} \int_0^\infty 
\left[ t \theta_3 \theta_3^{(1)} - t^2 \left( \theta_3^{(1)} \right)^2
+ 7 t^2 \theta_3 \theta_3^{(2)} + 6 t^3 \theta_3 \theta_3^{(3)}
-6 t^3 \theta_3^{(1)} \theta_3^{(2)}\right. \nonumber \\ & & \left.
 + t^4 \theta_3 \theta_3^{(4)}
- 4 t^4 \theta_3^{(1)} \theta_3^{(3)} + 3 t^4 \left( \theta_3^{(2)} \right)^2
\right]{\rm d}\mu_f\left(\frac{t}{A} \right),
\end{eqnarray}
and where the theta function and its derivatives are written for simplicity as
\begin{equation}
\theta_3 := \theta_3\left({\rm e}^{-t}\right) , \qquad
\theta_3^{(n)} := \frac{{\rm d}^n}{{\rm d}t^n}
\theta_3\left({\rm e}^{-t}\right), n\in \N.
\end{equation} 
\end{thm}

\begin{remark}
The same result can be written when $E_2$ is strictly increasing in the neighborhood of $A^*$, with the reverse condition $E_4(A^*)<0$, which simply exchanges the regimes of optimality.
\end{remark}

\begin{proof}
Since a transition point exists, a simple Taylor expansion of $E_2$ as $A\to A^*$ -- ensured by the fact that $E(\cdot,\Delta)\in C^\infty(\R_+^*)$ (see Remark \ref{regularity})-- reads
\begin{align*}
E_2(A)&=E_2(A^*)+\frac{d E_2}{d A}(A^*)(A-A^*)+o(A-A^*)\\
&=-\frac{d E_2}{d A}(A^*)(A^*-A)+o(A-A^*).
\end{align*}
By assumption, we know that $\frac{d E_2}{d A}(A^*)<0$ and let us write $b:=-\frac{d E_2}{d A}(A^*)>0$ for simplicity in such a way that, as $A\to A^*$,
\begin{equation}\label{E2}
E_2(A)=b(A^*-A)+o(A^*-A).
\end{equation}
The value of $\varepsilon$ which locally minimizes $E(A,\cdot)$ satisfies the following asymptotic stationarity condition we get from (\ref{ecreps}):
\begin{equation}\label{crit}
\frac{\partial E(A,e^\varepsilon)}{\partial \varepsilon}
= 2 E_2(A) \varepsilon + 4 E_4(A) \varepsilon^3 + O(\varepsilon^5) = 0 .    
\end{equation}
Using \eqref{E2}, this condition, for $A$ in the neighborhood of $A^*$, is therefore equivalent to the couple of equations
\begin{equation}
\varepsilon = 0 \qquad \textnormal{or} \qquad \varepsilon^2 =
\frac{b}{2 E_4(A^*)} \left( A - A^* \right) +o(A-A^*) .
\end{equation}  
We know that the first solution is dominant in a certain region
$A_0<A<A^*$ with $A_0\ge 0$.
Since $\varepsilon^2$ must be a real positive number and $E_4(A^*)>0$,
the second solution is
\begin{equation} \label{order}
\varepsilon= \sqrt{\frac{b}{2 E_4(A^*)}}
\left( A - A^* \right)^{1/2} +o(\left( A - A^* \right)^{1/2}) ,
\end{equation}  
dominant in the region $A^*<A<A_1$ for a certain $A_1>0$,
which completes the proof.
\end{proof}

\end{document}
